\documentclass{report}
\usepackage{amsmath,amssymb}
\setlength{\parindent}{0mm}
\setlength{\parskip}{1em}
\begin{document}
\begin{center}
\rule{6in}{1pt} \
{\large Tom Manteuffel \\
{\bf Range Decomposition: A low communication algorithm for solving PDEs on Massively Parallel Machines.}}

526 UCB \\ University of Colorado \\ Boulder \\ CO 80309
\\
{\tt tmanteuf@colorado.edu}\\
David J Appelhans\\
John Ruge\end{center}

The Range Decomposition (RD) algorithm uses nested iteration and adaptive
mesh refinement locally before performing a global communication step.
Only several such steps are observed to be necessary before reaching a
solution within a small multiple of discretization error. The target
application is peta- and exascale machines where traditional parallel
numerical PDE communication patterns stifle scalability. The RD algorithm
uses a partition of unity to equally distribute the error, and thus, the
work. The computational advantages of this approach are that the
decomposed problems can be solved using nested iteration and any
multigrid cycle type, in parallel, without any communication until the
partitioned solutions are summed. This offers potential advantages in the
paradigm of expensive communication but very cheap computation.
Performance models and numerical results demonstrate the enhanced
performance.


\end{document}
